% Options for packages loaded elsewhere
\PassOptionsToPackage{unicode}{hyperref}
\PassOptionsToPackage{hyphens}{url}
\documentclass[
]{article}
\usepackage{xcolor}
\usepackage[margin=1in]{geometry}
\usepackage{amsmath,amssymb}
\setcounter{secnumdepth}{5}
\usepackage{iftex}
\ifPDFTeX
  \usepackage[T1]{fontenc}
  \usepackage[utf8]{inputenc}
  \usepackage{textcomp} % provide euro and other symbols
\else % if luatex or xetex
  \usepackage{unicode-math} % this also loads fontspec
  \defaultfontfeatures{Scale=MatchLowercase}
  \defaultfontfeatures[\rmfamily]{Ligatures=TeX,Scale=1}
\fi
\usepackage{lmodern}
\ifPDFTeX\else
  % xetex/luatex font selection
\fi
% Use upquote if available, for straight quotes in verbatim environments
\IfFileExists{upquote.sty}{\usepackage{upquote}}{}
\IfFileExists{microtype.sty}{% use microtype if available
  \usepackage[]{microtype}
  \UseMicrotypeSet[protrusion]{basicmath} % disable protrusion for tt fonts
}{}
\makeatletter
\@ifundefined{KOMAClassName}{% if non-KOMA class
  \IfFileExists{parskip.sty}{%
    \usepackage{parskip}
  }{% else
    \setlength{\parindent}{0pt}
    \setlength{\parskip}{6pt plus 2pt minus 1pt}}
}{% if KOMA class
  \KOMAoptions{parskip=half}}
\makeatother
\usepackage{color}
\usepackage{fancyvrb}
\newcommand{\VerbBar}{|}
\newcommand{\VERB}{\Verb[commandchars=\\\{\}]}
\DefineVerbatimEnvironment{Highlighting}{Verbatim}{commandchars=\\\{\}}
% Add ',fontsize=\small' for more characters per line
\usepackage{framed}
\definecolor{shadecolor}{RGB}{248,248,248}
\newenvironment{Shaded}{\begin{snugshade}}{\end{snugshade}}
\newcommand{\AlertTok}[1]{\textcolor[rgb]{0.94,0.16,0.16}{#1}}
\newcommand{\AnnotationTok}[1]{\textcolor[rgb]{0.56,0.35,0.01}{\textbf{\textit{#1}}}}
\newcommand{\AttributeTok}[1]{\textcolor[rgb]{0.13,0.29,0.53}{#1}}
\newcommand{\BaseNTok}[1]{\textcolor[rgb]{0.00,0.00,0.81}{#1}}
\newcommand{\BuiltInTok}[1]{#1}
\newcommand{\CharTok}[1]{\textcolor[rgb]{0.31,0.60,0.02}{#1}}
\newcommand{\CommentTok}[1]{\textcolor[rgb]{0.56,0.35,0.01}{\textit{#1}}}
\newcommand{\CommentVarTok}[1]{\textcolor[rgb]{0.56,0.35,0.01}{\textbf{\textit{#1}}}}
\newcommand{\ConstantTok}[1]{\textcolor[rgb]{0.56,0.35,0.01}{#1}}
\newcommand{\ControlFlowTok}[1]{\textcolor[rgb]{0.13,0.29,0.53}{\textbf{#1}}}
\newcommand{\DataTypeTok}[1]{\textcolor[rgb]{0.13,0.29,0.53}{#1}}
\newcommand{\DecValTok}[1]{\textcolor[rgb]{0.00,0.00,0.81}{#1}}
\newcommand{\DocumentationTok}[1]{\textcolor[rgb]{0.56,0.35,0.01}{\textbf{\textit{#1}}}}
\newcommand{\ErrorTok}[1]{\textcolor[rgb]{0.64,0.00,0.00}{\textbf{#1}}}
\newcommand{\ExtensionTok}[1]{#1}
\newcommand{\FloatTok}[1]{\textcolor[rgb]{0.00,0.00,0.81}{#1}}
\newcommand{\FunctionTok}[1]{\textcolor[rgb]{0.13,0.29,0.53}{\textbf{#1}}}
\newcommand{\ImportTok}[1]{#1}
\newcommand{\InformationTok}[1]{\textcolor[rgb]{0.56,0.35,0.01}{\textbf{\textit{#1}}}}
\newcommand{\KeywordTok}[1]{\textcolor[rgb]{0.13,0.29,0.53}{\textbf{#1}}}
\newcommand{\NormalTok}[1]{#1}
\newcommand{\OperatorTok}[1]{\textcolor[rgb]{0.81,0.36,0.00}{\textbf{#1}}}
\newcommand{\OtherTok}[1]{\textcolor[rgb]{0.56,0.35,0.01}{#1}}
\newcommand{\PreprocessorTok}[1]{\textcolor[rgb]{0.56,0.35,0.01}{\textit{#1}}}
\newcommand{\RegionMarkerTok}[1]{#1}
\newcommand{\SpecialCharTok}[1]{\textcolor[rgb]{0.81,0.36,0.00}{\textbf{#1}}}
\newcommand{\SpecialStringTok}[1]{\textcolor[rgb]{0.31,0.60,0.02}{#1}}
\newcommand{\StringTok}[1]{\textcolor[rgb]{0.31,0.60,0.02}{#1}}
\newcommand{\VariableTok}[1]{\textcolor[rgb]{0.00,0.00,0.00}{#1}}
\newcommand{\VerbatimStringTok}[1]{\textcolor[rgb]{0.31,0.60,0.02}{#1}}
\newcommand{\WarningTok}[1]{\textcolor[rgb]{0.56,0.35,0.01}{\textbf{\textit{#1}}}}
\usepackage{longtable,booktabs,array}
\newcounter{none} % for unnumbered tables
\usepackage{calc} % for calculating minipage widths
% Correct order of tables after \paragraph or \subparagraph
\usepackage{etoolbox}
\makeatletter
\patchcmd\longtable{\par}{\if@noskipsec\mbox{}\fi\par}{}{}
\makeatother
% Allow footnotes in longtable head/foot
\IfFileExists{footnotehyper.sty}{\usepackage{footnotehyper}}{\usepackage{footnote}}
\makesavenoteenv{longtable}
\usepackage{graphicx}
\makeatletter
\newsavebox\pandoc@box
\newcommand*\pandocbounded[1]{% scales image to fit in text height/width
  \sbox\pandoc@box{#1}%
  \Gscale@div\@tempa{\textheight}{\dimexpr\ht\pandoc@box+\dp\pandoc@box\relax}%
  \Gscale@div\@tempb{\linewidth}{\wd\pandoc@box}%
  \ifdim\@tempb\p@<\@tempa\p@\let\@tempa\@tempb\fi% select the smaller of both
  \ifdim\@tempa\p@<\p@\scalebox{\@tempa}{\usebox\pandoc@box}%
  \else\usebox{\pandoc@box}%
  \fi%
}
% Set default figure placement to htbp
\def\fps@figure{htbp}
\makeatother
\setlength{\emergencystretch}{3em} % prevent overfull lines
\providecommand{\tightlist}{%
  \setlength{\itemsep}{0pt}\setlength{\parskip}{0pt}}
\usepackage{geometry}
\geometry{a4paper, margin=2.5cm}
\usepackage{fancyhdr}
\pagestyle{fancy}
\fancyhead[L]{Statistics Teaching System}
\fancyhead[R]{\thepage}
\fancyfoot[C]{}
\usepackage{booktabs}
\usepackage{longtable}
\usepackage{array}
\usepackage{hyperref}
\hypersetup{colorlinks=true, linkcolor=blue, urlcolor=blue}
\usepackage{xcolor}
\definecolor{codebg}{gray}{0.95}
\usepackage{listings}
\lstset{basicstyle=\ttfamily\small, breaklines=true, frame=single, backgroundcolor=\color{codebg}}
\usepackage{bookmark}
\IfFileExists{xurl.sty}{\usepackage{xurl}}{} % add URL line breaks if available
\urlstyle{same}
\hypersetup{
  pdftitle={Statistics Teaching System: Architecture and Operations Guide},
  pdfauthor={Yu-You Liou},
  hidelinks,
  pdfcreator={LaTeX via pandoc}}

\title{\textbf{Statistics Teaching System: Architecture and Operations
Guide}}
\author{Yu-You Liou}
\date{2026-06-07}

\begin{document}
\maketitle

{
\setcounter{tocdepth}{3}
\tableofcontents
}
\newpage

\section{System Overview}\label{system-overview}

This document describes the complete architecture and operational
procedures of the \textbf{Statistics Teaching System} (\texttt{stat}),
designed for the English Taught Program in International Business (ETP)
at Shih Chien University. The system supports two sequential
courses---\emph{Statistics (1)} and \emph{Statistics (2)}---and is built
on a core technology stack of R, Quarto, and R Markdown, integrating
website publication, in-class exercises, and automated grading into a
unified workflow.

\begin{quote}
\textbf{Course website:} \url{https://yyliou.github.io/stat}
\textbf{Instructor:} Yu-You Liou (\texttt{d10627008@ntu.edu.tw})
\end{quote}

\subsection{Course Structure}\label{course-structure}

\begin{longtable}[]{@{}
  >{\raggedright\arraybackslash}p{(\linewidth - 10\tabcolsep) * \real{0.1648}}
  >{\raggedright\arraybackslash}p{(\linewidth - 10\tabcolsep) * \real{0.1538}}
  >{\raggedright\arraybackslash}p{(\linewidth - 10\tabcolsep) * \real{0.1319}}
  >{\raggedright\arraybackslash}p{(\linewidth - 10\tabcolsep) * \real{0.3956}}
  >{\raggedright\arraybackslash}p{(\linewidth - 10\tabcolsep) * \real{0.0549}}
  >{\raggedright\arraybackslash}p{(\linewidth - 10\tabcolsep) * \real{0.0989}}@{}}
\caption{Course Identification and Scheduling}\tabularnewline
\toprule\noalign{}
\begin{minipage}[b]{\linewidth}\raggedright
Course
\end{minipage} & \begin{minipage}[b]{\linewidth}\raggedright
Code
\end{minipage} & \begin{minipage}[b]{\linewidth}\raggedright
Semester
\end{minipage} & \begin{minipage}[b]{\linewidth}\raggedright
Schedule
\end{minipage} & \begin{minipage}[b]{\linewidth}\raggedright
Room
\end{minipage} & \begin{minipage}[b]{\linewidth}\raggedright
Chapters
\end{minipage} \\
\midrule\noalign{}
\endfirsthead
\toprule\noalign{}
\begin{minipage}[b]{\linewidth}\raggedright
Course
\end{minipage} & \begin{minipage}[b]{\linewidth}\raggedright
Code
\end{minipage} & \begin{minipage}[b]{\linewidth}\raggedright
Semester
\end{minipage} & \begin{minipage}[b]{\linewidth}\raggedright
Schedule
\end{minipage} & \begin{minipage}[b]{\linewidth}\raggedright
Room
\end{minipage} & \begin{minipage}[b]{\linewidth}\raggedright
Chapters
\end{minipage} \\
\midrule\noalign{}
\endhead
\bottomrule\noalign{}
\endlastfoot
Statistics (1) & EIB-21E-01-A1 & Fall 2026 & Tuesday, Periods 8-9
(15:10-17:00) & A207 & Ch. 1-7 \\
Statistics (2) & EIB-21F-01-A1 & Spring 2026 & Thursday, Periods 6-7
(13:10-15:00) & L309 & Ch. 8-12 \\
\end{longtable}

Each session is structured as two consecutive hours: the first hour is a
\textbf{Lecture} covering theoretical foundations, and the second hour
is a \textbf{Lab} in which students apply concepts through in-class
exercises using the R programming language.

\subsection{Assessment Structure}\label{assessment-structure}

\begin{longtable}[]{@{}lll@{}}
\caption{Assessment Weighting and Components}\tabularnewline
\toprule\noalign{}
Component & Weight & Details \\
\midrule\noalign{}
\endfirsthead
\toprule\noalign{}
Component & Weight & Details \\
\midrule\noalign{}
\endhead
\bottomrule\noalign{}
\endlastfoot
In-Class Exercises & 36\% & 12 exercises x 3\% each \\
Midterm Examination & 32\% & \\
Final Examination & 32\% & \\
\end{longtable}

\newpage

\section{Directory Architecture}\label{directory-architecture}

\subsection{Root Directory Structure}\label{root-directory-structure}

The system root (\texttt{stat/}) is organized as follows:

\begin{verbatim}
stat/
|-- _quarto.yml          # Quarto website global configuration
|-- styles.css           # Custom CSS stylesheet (currently disabled in YAML)
|-- index.html           # Published homepage (rendered output)
|-- sylla.html           # Published syllabus landing page
|-- slide.html           # Published slide index
|-- prob.html            # Published problem set page
|-- r.html               # Published R guide
|-- site_libs/           # Quarto-generated site assets
|
|-- web/                 # Website page sources (qmd)
|   |-- web.Rproj        # RStudio project file
|   |-- index.qmd        # Homepage (course overview, exam regulations)
|   |-- sylla.qmd        # Syllabus landing page (links to PDF syllabi)
|   |-- slide.qmd        # Slide index page
|   |-- prob.qmd         # Problem set page (template, submission links)
|   `-- r.qmd            # R installation and usage guide
|
|-- ch2/ ... ch14/       # Per-chapter lecture slide directories
|   |-- chN.qmd          # English Reveal.js slide source
|   |-- chN-zh.qmd       # Traditional-Chinese parallel version
|   |-- bib.bib          # Chapter bibliography
|   `-- logo.svg         # Slide logo asset
|
|-- stat1-ps/            # Statistics (1) problem sets
|   |-- exN.qmd          # Exercise source (Reveal.js format)
|   `-- exN.html         # Compiled HTML for publication
|
|-- stat2-ps/            # Statistics (2) problem sets (same structure)
|
|-- stat1-ex/            # Statistics (1) student submissions
|   `-- ex1/ ... ex12/
|       `-- exN-[StudentID].pdf
|
|-- stat2-ex/            # Statistics (2) student submissions (same structure)
|
|-- stat1-mid/           # Statistics (1) midterm materials and grading outputs
|   |-- mid-score.xlsx   # Exam scoreboard (per-question scores)
|   |-- stat.Rmd         # Analytics report source
|   `-- stat.pdf         # Compiled analytics report
|-- stat1-fin/           # Statistics (1) final materials and grading outputs
|-- stat2-mid/           # Statistics (2) midterm materials and grading outputs
|-- stat2-fin/           # Statistics (2) final materials and grading outputs
|
|-- stu1/                # Statistics (1) roster and grading outputs
|   |-- [timestamp].xls  # Official roster exported from university system
|   |-- ex[N]/
|   |   |-- ex-score.xlsx
|   |   |-- stat.Rmd
|   |   `-- stat.pdf
|   `-- sem/             # Semester-aggregate outputs
|       |-- sem-score.xlsx
|       |-- stat.Rmd
|       `-- stat.pdf
|
|-- stu2/                # Statistics (2) roster and grading outputs (same structure)
|
|-- sylla/               # Syllabus source files
|   |-- stat1.Rmd / stat1.pdf
|   `-- stat2.Rmd / stat2.pdf
|
|-- !-Archive/           # Archive of retired materials
|
`-- prompt/              # AI-assisted prompt documents
    |-- slide-format.Rmd # Master specification for slide generation
    |-- prompt.Rmd       # Legacy slide generation prompt
    |-- ex-create.Rmd    # In-class exercise generation prompt
    |-- ex-review1.Rmd   # Exercise grading prompt for Statistics (1)
    |-- ex-review2.Rmd   # Exercise grading prompt for Statistics (2)
    |-- exam-review.Rmd  # Midterm/final exam auto-grading prompt
    |-- sem-review.Rmd   # Semester grade aggregation prompt
    |-- render.txt       # Batch render command for all chapter slides
    |-- bluman-step-by-step-statistics-8th-edition.pdf  # Textbook
    `-- stat-system-guide.Rmd   # This document
\end{verbatim}

\subsection{Design Principles}\label{design-principles}

The directory structure adheres to six organizational principles:

\begin{enumerate}
\def\labelenumi{\arabic{enumi}.}
\tightlist
\item
  \textbf{One directory per chapter.} Directories \texttt{ch2/} through
  \texttt{ch14/} are self-contained units (sources, bibliography, logo,
  rendered output), facilitating independent version control and
  selective rendering.
\item
  \textbf{Bilingual slide parity.} Every chapter maintains an English
  source (\texttt{chN.qmd}) and a structurally identical
  Traditional-Chinese version (\texttt{chN-zh.qmd}), cross-linked from
  the title slide.
\item
  \textbf{Separation of page sources and published output.} Website page
  sources live in \texttt{web/}; rendered HTML is published from the
  repository root for GitHub Pages.
\item
  \textbf{Separation of problem sets and submissions.}
  \texttt{stat1-ps/} (instructor-authored questions) and
  \texttt{stat1-ex/} (student submissions) maintain distinct
  responsibilities and access patterns.
\item
  \textbf{Co-location of rosters and grading outputs.} Each
  \texttt{stu1/} or \texttt{stu2/} directory holds both the official
  class roster and all automated grading artifacts, enabling
  single-location lookup.
\item
  \textbf{Centralized prompt management.} All AI-assisted documents,
  including the textbook PDF used as the generation source, are isolated
  within \texttt{prompt/}, decoupled from pedagogical content.
\end{enumerate}

\newpage

\section{Technology Stack}\label{technology-stack}

\subsection{Core Tools}\label{core-tools}

\begin{longtable}[]{@{}
  >{\raggedright\arraybackslash}p{(\linewidth - 4\tabcolsep) * \real{0.2000}}
  >{\raggedright\arraybackslash}p{(\linewidth - 4\tabcolsep) * \real{0.2000}}
  >{\raggedright\arraybackslash}p{(\linewidth - 4\tabcolsep) * \real{0.6000}}@{}}
\caption{Core Software Components}\tabularnewline
\toprule\noalign{}
\begin{minipage}[b]{\linewidth}\raggedright
Tool
\end{minipage} & \begin{minipage}[b]{\linewidth}\raggedright
Minimum Version
\end{minipage} & \begin{minipage}[b]{\linewidth}\raggedright
Purpose
\end{minipage} \\
\midrule\noalign{}
\endfirsthead
\toprule\noalign{}
\begin{minipage}[b]{\linewidth}\raggedright
Tool
\end{minipage} & \begin{minipage}[b]{\linewidth}\raggedright
Minimum Version
\end{minipage} & \begin{minipage}[b]{\linewidth}\raggedright
Purpose
\end{minipage} \\
\midrule\noalign{}
\endhead
\bottomrule\noalign{}
\endlastfoot
R & \textgreater= 4.2 & Statistical computation and document
generation \\
RStudio Desktop & Latest stable & Integrated development environment \\
Quarto & \textgreater= 1.3 & Website and slide publication \\
R Markdown & \textgreater= 2.20 & Assignment templates and grading
reports \\
TinyTeX & Current & LaTeX engine for PDF output \\
\end{longtable}

\subsection{Required R Packages}\label{required-r-packages}

The following packages must be installed prior to system operation:

\begin{Shaded}
\begin{Highlighting}[]
\CommentTok{\# Document generation}
\FunctionTok{install.packages}\NormalTok{(}\FunctionTok{c}\NormalTok{(}\StringTok{"rmarkdown"}\NormalTok{, }\StringTok{"knitr"}\NormalTok{, }\StringTok{"tinytex"}\NormalTok{))}
\NormalTok{tinytex}\SpecialCharTok{::}\FunctionTok{install\_tinytex}\NormalTok{()}

\CommentTok{\# Data manipulation and visualization}
\FunctionTok{install.packages}\NormalTok{(}\FunctionTok{c}\NormalTok{(}\StringTok{"tidyverse"}\NormalTok{, }\StringTok{"ggplot2"}\NormalTok{, }\StringTok{"patchwork"}\NormalTok{))}

\CommentTok{\# Spreadsheet I/O for grading reports}
\FunctionTok{install.packages}\NormalTok{(}\FunctionTok{c}\NormalTok{(}\StringTok{"readxl"}\NormalTok{, }\StringTok{"openxlsx"}\NormalTok{))}

\CommentTok{\# Academic integrity (UUID generation)}
\FunctionTok{install.packages}\NormalTok{(}\StringTok{"uuid"}\NormalTok{)}
\end{Highlighting}
\end{Shaded}

\subsection{Website Publication
Architecture}\label{website-publication-architecture}

The course website is built as a \textbf{Quarto Website} project,
governed by \texttt{\_quarto.yml}. Key configuration decisions are
summarized below.

\begin{longtable}[]{@{}ll@{}}
\caption{Quarto Website Configuration Summary}\tabularnewline
\toprule\noalign{}
Setting & Value \\
\midrule\noalign{}
\endfirsthead
\toprule\noalign{}
Setting & Value \\
\midrule\noalign{}
\endhead
\bottomrule\noalign{}
\endlastfoot
output-dir & . (HTML output at repository root for GitHub Pages) \\
Page sources & web/ (index, sylla, slide, prob, r) \\
Theme & Minty + brand (styles.css currently disabled) \\
Navigation bar & Syllabus -\textgreater{} Slides -\textgreater{} Problem
sets -\textgreater{} R \\
Slide format & Reveal.js \\
Math rendering & KaTeX \\
\end{longtable}

\newpage

\section{Lecture Slide Authoring
Workflow}\label{lecture-slide-authoring-workflow}

\subsection{File Format}\label{file-format}

Each chapter's slides reside in \texttt{chN/chN.qmd} (English) and
\texttt{chN/chN-zh.qmd} (Traditional Chinese). The YAML header is
\textbf{fixed} by the master specification
(\texttt{prompt/slide-format.Rmd}); only the subtitle and the
language-switch link target may vary:

\begin{Shaded}
\begin{Highlighting}[]
\PreprocessorTok{{-}{-}{-}}
\FunctionTok{title}\KeywordTok{:}\AttributeTok{ }\StringTok{"Statistics"}
\FunctionTok{subtitle}\KeywordTok{:}\AttributeTok{ }\StringTok{"Chapter N: [Title]"}
\FunctionTok{author}\KeywordTok{:}\AttributeTok{ }\StringTok{"Yu{-}You Liou"}
\FunctionTok{institute}\KeywordTok{:}\AttributeTok{ }\StringTok{"Shih Chien University"}
\FunctionTok{date}\KeywordTok{:}\AttributeTok{ }\StringTok{"2026{-}06{-}07"}
\FunctionTok{format}\KeywordTok{:}
\AttributeTok{  }\FunctionTok{revealjs}\KeywordTok{:}
\AttributeTok{    }\FunctionTok{html{-}math{-}method}\KeywordTok{:}\AttributeTok{ katex}
\AttributeTok{    }\FunctionTok{smaller}\KeywordTok{:}\AttributeTok{ }\CharTok{true}
\AttributeTok{    }\FunctionTok{theme}\KeywordTok{:}\AttributeTok{ serif}
\AttributeTok{    }\FunctionTok{footer}\KeywordTok{:}\AttributeTok{ }\StringTok{"Elementary Statistics: A Step by Step Approach"}
\AttributeTok{    }\FunctionTok{logo}\KeywordTok{:}\AttributeTok{ }\StringTok{"logo.svg"}
\AttributeTok{    }\FunctionTok{editor}\KeywordTok{:}\AttributeTok{ source}
\AttributeTok{    }\FunctionTok{slide{-}number}\KeywordTok{:}\AttributeTok{ }\CharTok{true}
\AttributeTok{    }\FunctionTok{scrollable}\KeywordTok{:}\AttributeTok{ }\CharTok{false}
\AttributeTok{    }\FunctionTok{transition}\KeywordTok{:}\AttributeTok{ }\StringTok{"none"}
\AttributeTok{    }\FunctionTok{menu}\KeywordTok{:}
\AttributeTok{      }\FunctionTok{useTextContentForMissingTitles}\KeywordTok{:}\AttributeTok{ }\CharTok{false}
\AttributeTok{      }\FunctionTok{hideMissingTitles}\KeywordTok{:}\AttributeTok{ }\CharTok{true}
\AttributeTok{      }\FunctionTok{side}\KeywordTok{:}\AttributeTok{ }\StringTok{\textquotesingle{}left\textquotesingle{}}
\AttributeTok{      }\FunctionTok{titleSelector}\KeywordTok{:}\AttributeTok{ }\StringTok{\textquotesingle{}h1\textquotesingle{}}
\PreprocessorTok{{-}{-}{-}}
\end{Highlighting}
\end{Shaded}

The header additionally includes an \texttt{include-after-body} script
and title-slide attributes that inject a language-switch link
(\texttt{chN-zh.html} in the English file; \texttt{chN.html} in the
Chinese file).

\subsection{Authoring Standards}\label{authoring-standards}

All slide content must conform to the master specification in
\texttt{prompt/slide-format.Rmd}, which is the single source of truth.
Key conventions:

\begin{itemize}
\tightlist
\item
  Each slide begins with a level-2 heading (\texttt{\#\#}).
\item
  Mathematical expressions use standard LaTeX syntax, e.g.,
  \texttt{\$\textbackslash{}bar\{x\}\ =\ \textbackslash{}frac\{\textbackslash{}sum\ x\}\{n\}\$}.
\item
  Definitions and key formulae are enclosed in Quarto callout blocks,
  which \textbf{must be perfectly balanced} (unbalanced callouts are the
  most common rendering failure):
\end{itemize}

\begin{Shaded}
\begin{Highlighting}[]
\NormalTok{::: \{.callout{-}note\}}
\NormalTok{**Definition** ...}
\NormalTok{:::}
\end{Highlighting}
\end{Shaded}

\begin{itemize}
\tightlist
\item
  Every textbook example provides two parallel solutions: (1) a complete
  algebraic derivation in LaTeX, and (2) an executable R implementation
  with displayed output.
\item
  All prose must be original; verbatim reproduction of textbook text is
  prohibited.
\item
  \textbf{Bilingual parity:} \texttt{chN.qmd} and \texttt{chN-zh.qmd}
  must be structurally identical, slide for slide.
\item
  Chapters \texttt{ch2/}--\texttt{ch5/} serve as the clean reference
  standard for layout and style.
\end{itemize}

\subsection{Generating a New Chapter}\label{generating-a-new-chapter}

\begin{verbatim}
Step 1. Create the chapter directory with shared assets:
        mkdir chN/ and copy bib.bib and logo.svg from an existing chapter.

Step 2. Provide prompt/slide-format.Rmd to an AI assistant and request
        chN.qmd and chN-zh.qmd for the target textbook chapter
        (textbook PDF available in prompt/). The agent must complete
        the validation checklist in the specification before delivery.

Step 3. Render both language versions (see prompt/render.txt for
        the batch command):
        quarto render chN/chN.qmd
        quarto render chN/chN-zh.qmd

Step 4. Update the chapter index in web/slide.qmd.
\end{verbatim}

\newpage

\section{Problem Set Management}\label{problem-set-management}

\subsection{Directory Mapping}\label{directory-mapping}

\begin{longtable}[]{@{}lll@{}}
\caption{Problem Set and Submission Directory Mapping}\tabularnewline
\toprule\noalign{}
Course & Problem Set Directory & Submission Directory \\
\midrule\noalign{}
\endfirsthead
\toprule\noalign{}
Course & Problem Set Directory & Submission Directory \\
\midrule\noalign{}
\endhead
\bottomrule\noalign{}
\endlastfoot
Statistics (1) & stat1-ps/ & stat1-ex/ \\
Statistics (2) & stat2-ps/ & stat2-ex/ \\
\end{longtable}

\subsection{Problem Set Format}\label{problem-set-format}

Problem sets are authored as Reveal.js presentations
(\texttt{stat1-ps/exN.qmd}), with one question per slide. A
representative header:

\begin{Shaded}
\begin{Highlighting}[]
\PreprocessorTok{{-}{-}{-}}
\FunctionTok{title}\KeywordTok{:}\AttributeTok{ }\StringTok{"Statistics 1"}
\FunctionTok{subtitle}\KeywordTok{:}\AttributeTok{ }\StringTok{"In{-}Class Exercise N"}
\FunctionTok{author}\KeywordTok{:}\AttributeTok{ }\StringTok{"Yu{-}You Liou"}
\FunctionTok{date}\KeywordTok{:}\AttributeTok{ today}
\FunctionTok{institute}\KeywordTok{:}\AttributeTok{ }\StringTok{"Shih Chien University"}
\FunctionTok{format}\KeywordTok{:}
\AttributeTok{  }\FunctionTok{revealjs}\KeywordTok{:}
\AttributeTok{    }\FunctionTok{html{-}math{-}method}\KeywordTok{:}\AttributeTok{ katex}
\PreprocessorTok{{-}{-}{-}}
\end{Highlighting}
\end{Shaded}

\subsection{Submission and Grade Distribution
Channels}\label{submission-and-grade-distribution-channels}

Submission and grade lookup are handled through shared \textbf{pCloud}
links published on the problem set page (\texttt{web/prob.qmd}): one
upload link per course (students place the PDF in the matching
\texttt{exN} folder, which syncs to \texttt{stat{[}N{]}-ex/exN/}), and
one read-only link per course for checking grades.

\subsection{Submission Requirements}\label{submission-requirements}

Students must use the official template below and upload to the course
pCloud folder \texttt{exN/} (synced to \texttt{stat1-ex/exN/}) using the
exact naming convention \texttt{exN-{[}StudentID{]}.pdf}.

\begin{Shaded}
\begin{Highlighting}[]
\SpecialCharTok{{-}{-}{-}}
\NormalTok{title}\SpecialCharTok{:} \StringTok{"Ex. N"}
\NormalTok{author}\SpecialCharTok{:} \StringTok{"StudentID"}
\NormalTok{date}\SpecialCharTok{:} \StringTok{"2026{-}06{-}07"}
\NormalTok{output}\SpecialCharTok{:}\NormalTok{ pdf\_document}
\SpecialCharTok{{-}{-}{-}}

\StringTok{\textasciigrave{}}\AttributeTok{ }\StringTok{\textasciigrave{}} \StringTok{\textasciigrave{}}\AttributeTok{\{r setup, include=FALSE\}}
\AttributeTok{knitr::opts\_chunk$set(echo = TRUE)}
\StringTok{\textasciigrave{}} \StringTok{\textasciigrave{}}\AttributeTok{ }\StringTok{\textasciigrave{}}

\StringTok{\textasciigrave{}}\AttributeTok{ }\StringTok{\textasciigrave{}} \StringTok{\textasciigrave{}}\AttributeTok{\{r\}}
\AttributeTok{print(uuid::UUIDgenerate())}
\StringTok{\textasciigrave{}} \StringTok{\textasciigrave{}}\AttributeTok{ }\StringTok{\textasciigrave{}}
\end{Highlighting}
\end{Shaded}

\begin{quote}
\textbf{Every submission must include the UUID generation code block.
This serves as the primary mechanism for academic integrity
verification.}
\end{quote}

\newpage

\section{Automated Grading System}\label{automated-grading-system}

The system comprises three AI-driven grading pipelines: (1) in-class
exercise grading (\texttt{ex-review1.Rmd} / \texttt{ex-review2.Rmd}),
operating on PDF submissions; (2) examination grading
(\texttt{exam-review.Rmd}), operating on Google Form responses; and (3)
semester aggregation (\texttt{sem-review.Rmd}), which consolidates all
components into the final grade.

\subsection{Exercise Grading: Process
Overview}\label{exercise-grading-process-overview}

The grading pipeline consists of four sequential stages:

\begin{verbatim}
STAGE 1 | Initialization and Roster Reconciliation
         Read stu[N]/[timestamp].xls -> construct student roster
         Cross-reference stat[N]-ex/exN/ -> flag absent students
         as Missing (Score = 0)

STAGE 2 | Compliance Filtering (sequential; stops at first violation)
         (i)  Submission timestamp > [deadline] (UTC+8)
              -> Late Submission (Score = 0)
         (ii) File is not PDF or does not match exN-[StudentID].pdf
              -> Incorrect File Name/Format (Score = 0)
         (iii) YAML missing title/author/date, or UUID chunk absent,
               or author field does not match submitting student ID
              -> Incorrect Internal Format (Score = 0)
         (iv) UUID value identical to another submission, or
              invalid/forged UUID
              -> Plagiarism (Score = 0; all implicated students)

STAGE 3 | Academic Content Grading
         Total: 100 points, allocated equally across N questions
         Each question scored against stat[N]-ps/exN.qmd answer key
         Incorrect items noted as: [QN] is incorrect

STAGE 4 | Output Generation
         stu[N]/exN/ex-score.xlsx  <- Scoreboard (4 columns)
         stu[N]/exN/stat.Rmd       <- Analytics report source
         stu[N]/exN/stat.pdf       <- Compiled analytics report
\end{verbatim}

\subsection{\texorpdfstring{Scoreboard Format
(\texttt{ex-score.xlsx})}{Scoreboard Format (ex-score.xlsx)}}\label{scoreboard-format-ex-score.xlsx}

\begin{longtable}[]{@{}
  >{\raggedright\arraybackslash}p{(\linewidth - 4\tabcolsep) * \real{0.1220}}
  >{\raggedright\arraybackslash}p{(\linewidth - 4\tabcolsep) * \real{0.1951}}
  >{\raggedright\arraybackslash}p{(\linewidth - 4\tabcolsep) * \real{0.6829}}@{}}
\caption{Scoreboard Column Specification}\tabularnewline
\toprule\noalign{}
\begin{minipage}[b]{\linewidth}\raggedright
Column
\end{minipage} & \begin{minipage}[b]{\linewidth}\raggedright
Type
\end{minipage} & \begin{minipage}[b]{\linewidth}\raggedright
Description
\end{minipage} \\
\midrule\noalign{}
\endfirsthead
\toprule\noalign{}
\begin{minipage}[b]{\linewidth}\raggedright
Column
\end{minipage} & \begin{minipage}[b]{\linewidth}\raggedright
Type
\end{minipage} & \begin{minipage}[b]{\linewidth}\raggedright
Description
\end{minipage} \\
\midrule\noalign{}
\endhead
\bottomrule\noalign{}
\endlastfoot
StudentID & Character & University student identification number \\
Name & Character & Student full name as listed in official roster \\
Score & Integer (0-100) & Final score awarded for the submission \\
Note & Character & Compliance violation code or incorrect question
numbers \\
\end{longtable}

\subsection{\texorpdfstring{Analytics Report
(\texttt{stat.Rmd})}{Analytics Report (stat.Rmd)}}\label{analytics-report-stat.rmd}

The analytics report reads \texttt{ex-score.xlsx} and produces three
outputs:

\begin{enumerate}
\def\labelenumi{\arabic{enumi}.}
\tightlist
\item
  \textbf{Descriptive statistics table:} minimum, \(Q_1\), median,
  \(Q_3\), maximum, and mean of the score distribution.
\item
  \textbf{Boxplot:} a publication-quality visualization of the score
  distribution rendered via \texttt{ggplot2}.
\item
  \textbf{Submission summary table:} counts for Total, Missing, Late
  Submission, Incorrect Format, and Plagiarism.
\end{enumerate}

\subsection{Grading Note Codes}\label{grading-note-codes}

\begin{longtable}[]{@{}
  >{\raggedright\arraybackslash}p{(\linewidth - 4\tabcolsep) * \real{0.2432}}
  >{\raggedright\arraybackslash}p{(\linewidth - 4\tabcolsep) * \real{0.5946}}
  >{\raggedright\arraybackslash}p{(\linewidth - 4\tabcolsep) * \real{0.1622}}@{}}
\caption{Grading Note Code Reference}\tabularnewline
\toprule\noalign{}
\begin{minipage}[b]{\linewidth}\raggedright
Code
\end{minipage} & \begin{minipage}[b]{\linewidth}\raggedright
Description
\end{minipage} & \begin{minipage}[b]{\linewidth}\raggedright
Score
\end{minipage} \\
\midrule\noalign{}
\endfirsthead
\toprule\noalign{}
\begin{minipage}[b]{\linewidth}\raggedright
Code
\end{minipage} & \begin{minipage}[b]{\linewidth}\raggedright
Description
\end{minipage} & \begin{minipage}[b]{\linewidth}\raggedright
Score
\end{minipage} \\
\midrule\noalign{}
\endhead
\bottomrule\noalign{}
\endlastfoot
Missing & No file submitted by the deadline & 0 \\
Late Submission & File submitted after the official deadline & 0 \\
Incorrect File Name/Format & File is not PDF or does not follow naming
convention exN-{[}ID{]}.pdf & 0 \\
Incorrect Internal Format & YAML header incomplete, author ID mismatch,
or UUID chunk absent & 0 \\
Plagiarism & UUID value matches another submission, or is invalid/forged
& 0 \\
{[}QN{]} is incorrect & Question N answer is incorrect or methodology is
flawed & Partial deduction \\
\end{longtable}

\subsection{\texorpdfstring{Examination Grading Pipeline
(\texttt{exam-review.Rmd})}{Examination Grading Pipeline (exam-review.Rmd)}}\label{examination-grading-pipeline-exam-review.rmd}

Examinations are submitted online through a Google Form whose responses
populate a Google Sheet with the fixed schema:

\begin{verbatim}
time | ip | user | id | name | q1 | q1ans | q2 | q2ans | ... | qN | qNans
\end{verbatim}

where \texttt{qk} holds the student's R code for Question \(k\) and
\texttt{qkans} its reported output. The pipeline takes three parameters
--- \texttt{{[}sheet{]}} (the Google Sheet URL),
\texttt{{[}questions{]}} (the official question set with answers), and
\texttt{{[}exam{]}} (\texttt{mid} or \texttt{fin}) --- plus the campus
IP range constant \texttt{{[}campus-ip{]}}.

\begin{verbatim}
STAGE 1 | Data Acquisition and Roster Reconciliation
         Download [sheet] as CSV; infer N from column pairs;
         keep only the latest row per student ID.
         Roster students with no submission -> Missing (Score = 0)

STAGE 2 | Integrity Filtering (sequential; stops at first violation)
         (i)  ip outside [campus-ip]
              -> Invalid IP (off-campus) (Score = 0)
         (ii) user value shared across different student IDs
              -> Duplicate device (Score = 0; all implicated students)

STAGE 3 | Per-Question Grading with Partial Credit
         Total: 100 points, allocated equally across N questions.
         Each question's score is recorded individually.
         Partial-credit rubric:
           Full -> correct method and matching answer
           2/3  -> correct method/code, minor slip in reported answer
           1/3  -> correct method, mis-transcribed input values
           0    -> wrong method, irrelevant code, or blank

STAGE 4 | Output Generation ([dir] = stat[N]-mid/ or stat[N]-fin/)
         [dir]/[exam]-score.xlsx  <- Scoreboard (ID, Name, Q1..QN,
                                     Score, Note)
         [dir]/stat.Rmd           <- Analytics report source
         [dir]/stat.pdf           <- Compiled analytics report
\end{verbatim}

Compared to exercise grading, the exam analytics report adds a
\textbf{per-question average-score bar chart}, and the submissions
summary counts Invalid IP and Duplicate Device cases instead of
Late/Plagiarism.

\subsection{\texorpdfstring{Semester Grade Aggregation
(\texttt{sem-review.Rmd})}{Semester Grade Aggregation (sem-review.Rmd)}}\label{semester-grade-aggregation-sem-review.rmd}

Run once per semester after all components are graded. Parameters:
\texttt{{[}sem{]}} (1 or 2) and the component weights
\texttt{{[}w-mid{]}}, \texttt{{[}w-fin{]}}, \texttt{{[}w-ex{]}} (must
sum to 100; current syllabus values 32/32/36).

\begin{verbatim}
INPUT    stat[sem]-mid/mid-score.xlsx        (midterm)
         stat[sem]-fin/fin-score.xlsx        (final)
         stu[sem]/ex1..ex12/ex-score.xlsx    (12 exercises)

COMPUTE  ExAvg = mean(Ex1..Ex12)
         Total = w-mid/100 * Mid + w-fin/100 * Fin + w-ex/100 * ExAvg
         (rounded to 2 decimals)

OUTPUT   stu[sem]/sem/sem-score.xlsx
           Columns: StudentID, Name, Ex1..Ex12, ExAvg, Mid, Fin,
                    Total, Note
         stu[sem]/sem/stat.Rmd + stat.pdf
           Descriptive statistics of Total; boxplot; component-mean
           bar chart; histogram (bin width 10); pass/fail summary
           (Total >= 60); weights recap
\end{verbatim}

Robustness rules: the roster is the master index (no student silently
dropped); a missing component file or absent student contributes 0 and
is flagged in \texttt{Note}; zero-score trigger reasons are carried over
in compact form (e.g., \texttt{Mid:\ Invalid\ IP}); 3 randomly sampled
Totals are re-computed by hand for verification.

\newpage

\section{Academic Integrity
Mechanisms}\label{academic-integrity-mechanisms}

\subsection{UUID Verification}\label{uuid-verification}

Every submission must execute and display output from the following R
code block:

\begin{Shaded}
\begin{Highlighting}[]
\FunctionTok{print}\NormalTok{(uuid}\SpecialCharTok{::}\FunctionTok{UUIDgenerate}\NormalTok{())}
\end{Highlighting}
\end{Shaded}

\texttt{uuid::UUIDgenerate()} produces a Version 4 UUID (Universally
Unique Identifier) on each invocation, e.g.,
\texttt{6d2b9fda-2d9b-48f3-84bc-f82eabc39b4d}. The grading system
compares this value across all submissions for a given exercise. If two
or more files share an identical UUID, or if a file's UUID is invalid or
forged, all implicated students receive a score of zero with the note
\texttt{Plagiarism}.

\textbf{Important implementation notes:}

\begin{itemize}
\tightlist
\item
  The compliance filter is applied sequentially per file. A file that
  fails Stage 2 check (iii) (Incorrect Internal Format) is not evaluated
  for plagiarism within its own record; however, its UUID remains in the
  comparison pool and may trigger a Plagiarism flag in other students'
  records.
\item
  The \texttt{author} field in the YAML header must contain the
  submitting student's own ID. Use of another student's ID constitutes a
  format violation independent of the plagiarism check.
\end{itemize}

\subsection{Examination Regulations}\label{examination-regulations}

In alignment with the course's paperless policy, examinations are
conducted online (a sample examination interface is linked from the
course homepage). Examination materials are maintained in
\texttt{stat1-mid/}, \texttt{stat1-fin/}, \texttt{stat2-mid/}, and
\texttt{stat2-fin/}.

\begin{longtable}[]{@{}
  >{\raggedright\arraybackslash}p{(\linewidth - 2\tabcolsep) * \real{0.2321}}
  >{\raggedright\arraybackslash}p{(\linewidth - 2\tabcolsep) * \real{0.7679}}@{}}
\caption{Examination Conduct Regulations}\tabularnewline
\toprule\noalign{}
\begin{minipage}[b]{\linewidth}\raggedright
Regulation
\end{minipage} & \begin{minipage}[b]{\linewidth}\raggedright
Policy
\end{minipage} \\
\midrule\noalign{}
\endfirsthead
\toprule\noalign{}
\begin{minipage}[b]{\linewidth}\raggedright
Regulation
\end{minipage} & \begin{minipage}[b]{\linewidth}\raggedright
Policy
\end{minipage} \\
\midrule\noalign{}
\endhead
\bottomrule\noalign{}
\endlastfoot
Mandatory attendance & Roll call 5 minutes before start; late arrival or
absence results in score = 0 \\
Network restriction & System tracks IP addresses; campus Wi-Fi required;
non-campus IP results in score = 0 \\
Authorized resources & Open-book; reference materials and AI tools are
permitted as auxiliary aids \\
Media prohibition & Photography, video/audio recording, and
screenshotting result in score = 0 \\
Communication prohibition & Interpersonal communication and disruptive
behavior result in score = 0 \\
Copy-paste disabled & Copy-paste functionality is disabled by the
examination system \\
\end{longtable}

\subsection{AI Use Policy}\label{ai-use-policy}

Students may use AI tools (e.g., ChatGPT, Gemini) as learning aids.
However, students must be capable of explaining and defending any code
or analysis they submit. Uncritical submission of AI-generated content
without demonstrated understanding constitutes academic dishonesty and
is subject to the same consequences as plagiarism.

\newpage

\section{AI Prompt Management}\label{ai-prompt-management}

The \texttt{prompt/} directory contains structured prompt documents for
AI-assisted course material generation and grading.

\begin{longtable}[]{@{}
  >{\raggedright\arraybackslash}p{(\linewidth - 2\tabcolsep) * \real{0.2093}}
  >{\raggedright\arraybackslash}p{(\linewidth - 2\tabcolsep) * \real{0.7907}}@{}}
\caption{Prompt Directory File Inventory}\tabularnewline
\toprule\noalign{}
\begin{minipage}[b]{\linewidth}\raggedright
File
\end{minipage} & \begin{minipage}[b]{\linewidth}\raggedright
Purpose
\end{minipage} \\
\midrule\noalign{}
\endfirsthead
\toprule\noalign{}
\begin{minipage}[b]{\linewidth}\raggedright
File
\end{minipage} & \begin{minipage}[b]{\linewidth}\raggedright
Purpose
\end{minipage} \\
\midrule\noalign{}
\endhead
\bottomrule\noalign{}
\endlastfoot
slide-format.Rmd & Master specification for AI-assisted slide generation
(single source of truth; supersedes prompt.Rmd) \\
prompt.Rmd & Legacy slide generation prompt (retained for reference) \\
ex-create.Rmd & In-class exercise generation prompt (textbook-scoped,
template-bound) \\
ex-review1.Rmd & Automated grading prompt for Statistics (1)
exercises \\
ex-review2.Rmd & Automated grading prompt for Statistics (2)
exercises \\
exam-review.Rmd & Midterm/final examination auto-grading prompt (Google
Form pipeline) \\
sem-review.Rmd & Semester grade aggregation prompt (weighted total of
all components) \\
render.txt & Batch quarto render command for all chapter slides \\
bluman-\ldots-8th-edition.pdf & Source textbook PDF used for slide
generation \\
stat-system-guide.Rmd & This technical document \\
\end{longtable}

\subsection{\texorpdfstring{Slide Generation Specification
(\texttt{slide-format.Rmd})}{Slide Generation Specification (slide-format.Rmd)}}\label{slide-generation-specification-slide-format.rmd}

\texttt{slide-format.Rmd} is a complete, self-contained specification
pasted to the AI agent together with the instruction \emph{``Generate
\texttt{chX.qmd} (and \texttt{chX-zh.qmd}) for Chapter X following this
specification exactly.''} Key directives:

\begin{itemize}
\tightlist
\item
  \textbf{Fixed YAML front matter:} Only the subtitle and the
  language-switch link may vary; everything else is reproduced verbatim.
\item
  \textbf{Structural correctness first:} Decks must render without
  error; every callout block must be perfectly balanced (the most common
  failure mode).
\item
  \textbf{Exhaustive extraction:} Every example in the target chapter
  must be incorporated; summarization or omission is not permitted.
\item
  \textbf{Dual-solution requirement:} Each example must present (1) a
  full algebraic derivation in LaTeX and (2) executable R code with
  displayed output.
\item
  \textbf{Paraphrasing mandate:} All prose must be rewritten in original
  language; verbatim textbook reproduction is prohibited.
\item
  \textbf{Bilingual parity:} Each chapter ships as a structurally
  identical English/Traditional-Chinese pair.
\item
  \textbf{Reference standard:} Chapters \texttt{ch2}--\texttt{ch5} are
  the clean style reference.
\item
  \textbf{Validation checklist:} The agent must run the specification's
  built-in checklist before returning a file.
\end{itemize}

\subsection{\texorpdfstring{Exercise Generation Prompt
(\texttt{ex-create.Rmd})}{Exercise Generation Prompt (ex-create.Rmd)}}\label{exercise-generation-prompt-ex-create.rmd}

Generates a new in-class exercise deck from the textbook. Parameters:
\texttt{{[}no{]}} (exercise number), \texttt{{[}range{]}} (textbook
coverage, Chapter.Section to Chapter.Section), and \texttt{{[}n{]}}
(number of questions). Key directives:

\begin{itemize}
\tightlist
\item
  Questions are scoped strictly to \texttt{{[}range{]}}, balanced across
  its sections, and ordered easiest to hardest.
\item
  Every question must be solvable in a few lines of R, with all required
  data given inline.
\item
  All wording and numbers must be original (no verbatim textbook reuse).
\item
  Output replicates \texttt{stat1-ps/ex1.qmd} exactly (one question per
  Reveal.js slide, callout-warning blocks) and is saved as
  \texttt{stat{[}N{]}-ps/ex{[}no{]}.qmd}, then rendered to HTML with a
  built-in verification checklist.
\end{itemize}

\subsection{\texorpdfstring{Exercise Grading Prompts
(\texttt{ex-review1.Rmd},
\texttt{ex-review2.Rmd})}{Exercise Grading Prompts (ex-review1.Rmd, ex-review2.Rmd)}}\label{exercise-grading-prompts-ex-review1.rmd-ex-review2.rmd}

Key components of the grading prompts:

\begin{itemize}
\tightlist
\item
  Definition of environment variables: \texttt{{[}deadline{]}} (UTC+8),
  \texttt{{[}no{]}} (exercise number), and root directory structure.
\item
  Specification of all four grading stages with precise conditional
  logic.
\item
  Output format requirements: scoreboard column names, Rmd content
  sections, and PDF compilation.
\end{itemize}

The two files are identical except for the directory targets
(\texttt{stu1}/\texttt{stat1-ps}/\texttt{stat1-ex}
vs.~\texttt{stu2}/\texttt{stat2-ps}/\texttt{stat2-ex}).

\subsection{\texorpdfstring{Examination Grading Prompt
(\texttt{exam-review.Rmd})}{Examination Grading Prompt (exam-review.Rmd)}}\label{examination-grading-prompt-exam-review.rmd}

Parameters: \texttt{{[}sheet{]}} (Google Sheet URL of form responses),
\texttt{{[}questions{]}} (official question set with answers),
\texttt{{[}exam{]}} (\texttt{mid}/\texttt{fin}), and the constant
\texttt{{[}campus-ip{]}}. Implements the four-stage examination pipeline
described in the Automated Grading System section: CSV acquisition and
roster reconciliation, integrity filtering (campus IP, duplicate
device), per-question grading with the full/2-3/1-3/0 partial-credit
rubric, and output to \texttt{stat{[}N{]}-mid/} or
\texttt{stat{[}N{]}-fin/}.

\subsection{\texorpdfstring{Semester Aggregation Prompt
(\texttt{sem-review.Rmd})}{Semester Aggregation Prompt (sem-review.Rmd)}}\label{semester-aggregation-prompt-sem-review.rmd}

Parameters: \texttt{{[}sem{]}} (1/2) and weights \texttt{{[}w-mid{]}},
\texttt{{[}w-fin{]}}, \texttt{{[}w-ex{]}} (sum = 100). Consolidates the
midterm, final, and 12 exercise scoreboards into
\texttt{stu{[}sem{]}/sem/sem-score.xlsx} with the weighted Total,
carries over zero-score reasons into the Note column, and produces the
semester analytics report (\texttt{stat.Rmd}/\texttt{stat.pdf}).

\newpage

\section{Operations Guide}\label{operations-guide}

\subsection{New Semester
Initialization}\label{new-semester-initialization}

\begin{verbatim}
1. Export official class roster from the university administrative system.
   Save to stu1/ (Statistics 1) or stu2/ (Statistics 2).

2. Update sylla/stat1.Rmd and sylla/stat2.Rmd with the current
   semester's schedule and policies. Knit to PDF.

3. Verify _quarto.yml navigation links and semester metadata in
   web/index.qmd and web/sylla.qmd.

4. Create or empty the pCloud submission folders (ex1-ex12) and
   confirm the upload/grade-lookup links in web/prob.qmd.

5. Rebuild the website pages and republish.
\end{verbatim}

\subsection{Weekly Problem Set
Publication}\label{weekly-problem-set-publication}

\begin{verbatim}
1. Create stat1-ps/exN.qmd (or stat2-ps/exN.qmd) in Reveal.js format,
   either manually or by executing the AI generation prompt
   (prompt/ex-create.Rmd) with [no], [range], and [n].
   Each slide corresponds to one question.

2. Render the problem set:
   quarto render stat1-ps/exN.qmd

3. Update the problem set index in web/prob.qmd.
\end{verbatim}

\subsection{Weekly Submission Collection and
Grading}\label{weekly-submission-collection-and-grading}

\begin{verbatim}
1. Students upload exN-[StudentID].pdf to the course pCloud folder,
   which syncs to stat1-ex/exN/.

2. Execute the AI grading prompt (prompt/ex-review1.Rmd) with
   [no] set to the exercise number and [deadline] set to the
   submission deadline (UTC+8).

3. Verify generated scoreboard: stu1/exN/ex-score.xlsx

4. Verify analytics report: stu1/exN/stat.pdf

5. Grades are published to students via the read-only pCloud link.
\end{verbatim}

\subsection{Examination Grading}\label{examination-grading}

\begin{verbatim}
1. Conduct the exam online; responses collect in the Google Sheet.

2. Execute the AI exam grading prompt (prompt/exam-review.Rmd) with
   [sheet] set to the response sheet URL, [questions] set to the
   official question set (with answers), and [exam] = mid or fin.

3. Verify scoreboard: stat[N]-mid/mid-score.xlsx (or stat[N]-fin/
   fin-score.xlsx) — per-question scores plus zero-score reasons.

4. Verify analytics report: stat.pdf in the same directory.
\end{verbatim}

\subsection{End-of-Semester Grade
Aggregation}\label{end-of-semester-grade-aggregation}

\begin{verbatim}
1. Confirm all 12 exercise scoreboards and both exam scoreboards exist.

2. Execute the AI aggregation prompt (prompt/sem-review.Rmd) with
   [sem] = 1 or 2 and the syllabus weights ([w-mid]/[w-fin]/[w-ex]
   = 32/32/36).

3. Verify stu[sem]/sem/sem-score.xlsx (row count = roster size) and
   the semester analytics report stat.pdf.

4. Submit final grades to the university system.
\end{verbatim}

\subsection{Generating a New Lecture
Chapter}\label{generating-a-new-lecture-chapter}

\begin{verbatim}
1. Create chN/ and copy bib.bib and logo.svg from an existing chapter.

2. Generate chN.qmd and chN-zh.qmd using prompt/slide-format.Rmd
   with an AI assistant, referencing the textbook PDF in prompt/.

3. Review, edit, and finalize both files in RStudio.

4. Render (batch command in prompt/render.txt):
   quarto render 'ch*/ch*.qmd'

5. Add the chapter link to web/slide.qmd.
\end{verbatim}

\subsection{Full-Site Publication (GitHub
Pages)}\label{full-site-publication-github-pages}

\begin{Shaded}
\begin{Highlighting}[]
\CommentTok{\# Render updated pages/slides, e.g.}
\ExtensionTok{quarto}\NormalTok{ render }\StringTok{\textquotesingle{}ch*/ch*.qmd\textquotesingle{}}

\CommentTok{\# Commit and push to trigger GitHub Pages deployment}
\FunctionTok{git}\NormalTok{ add .}
\FunctionTok{git}\NormalTok{ commit }\AttributeTok{{-}m} \StringTok{"Update: Week N — [brief description]"}
\FunctionTok{git}\NormalTok{ push origin main}
\end{Highlighting}
\end{Shaded}

\subsection{Troubleshooting Reference}\label{troubleshooting-reference}

\begin{longtable}[]{@{}
  >{\raggedright\arraybackslash}p{(\linewidth - 4\tabcolsep) * \real{0.2500}}
  >{\raggedright\arraybackslash}p{(\linewidth - 4\tabcolsep) * \real{0.3600}}
  >{\raggedright\arraybackslash}p{(\linewidth - 4\tabcolsep) * \real{0.3900}}@{}}
\caption{Common Issues and Resolutions}\tabularnewline
\toprule\noalign{}
\begin{minipage}[b]{\linewidth}\raggedright
Symptom
\end{minipage} & \begin{minipage}[b]{\linewidth}\raggedright
Probable Cause
\end{minipage} & \begin{minipage}[b]{\linewidth}\raggedright
Resolution
\end{minipage} \\
\midrule\noalign{}
\endfirsthead
\toprule\noalign{}
\begin{minipage}[b]{\linewidth}\raggedright
Symptom
\end{minipage} & \begin{minipage}[b]{\linewidth}\raggedright
Probable Cause
\end{minipage} & \begin{minipage}[b]{\linewidth}\raggedright
Resolution
\end{minipage} \\
\midrule\noalign{}
\endhead
\bottomrule\noalign{}
\endlastfoot
PDF output fails & TinyTeX not installed or incomplete & Run:
tinytex::install\_tinytex() \\
UUID package not found & uuid package not installed & Run:
install.packages(`uuid') \\
Quarto render error & Quarto version outdated & Update Quarto to latest
stable release \\
Roster file read error & Legacy .xls format compatibility & Use readxl
package or resave as .xlsx \\
Slide math not rendering & KaTeX configuration missing & Set
html-math-method: katex in YAML \\
\end{longtable}

\newpage

\section{Course Chapter Reference}\label{course-chapter-reference}

\subsection{Statistics (1): Chapter
Schedule}\label{statistics-1-chapter-schedule}

\begin{longtable}[]{@{}
  >{\raggedright\arraybackslash}p{(\linewidth - 6\tabcolsep) * \real{0.0833}}
  >{\raggedright\arraybackslash}p{(\linewidth - 6\tabcolsep) * \real{0.1111}}
  >{\raggedright\arraybackslash}p{(\linewidth - 6\tabcolsep) * \real{0.6806}}
  >{\raggedright\arraybackslash}p{(\linewidth - 6\tabcolsep) * \real{0.1250}}@{}}
\caption{Statistics (1) Weekly Schedule and Exercise
Mapping}\tabularnewline
\toprule\noalign{}
\begin{minipage}[b]{\linewidth}\raggedright
Week
\end{minipage} & \begin{minipage}[b]{\linewidth}\raggedright
Chapter
\end{minipage} & \begin{minipage}[b]{\linewidth}\raggedright
Topic
\end{minipage} & \begin{minipage}[b]{\linewidth}\raggedright
Exercise
\end{minipage} \\
\midrule\noalign{}
\endfirsthead
\toprule\noalign{}
\begin{minipage}[b]{\linewidth}\raggedright
Week
\end{minipage} & \begin{minipage}[b]{\linewidth}\raggedright
Chapter
\end{minipage} & \begin{minipage}[b]{\linewidth}\raggedright
Topic
\end{minipage} & \begin{minipage}[b]{\linewidth}\raggedright
Exercise
\end{minipage} \\
\midrule\noalign{}
\endhead
\bottomrule\noalign{}
\endlastfoot
1 & --- & Syllabus and Class Regulations & --- \\
2 & Ch. 1 & The Nature of Probability and Statistics & ex1 \\
3 & Ch. 2 & Frequency Distributions and Graphs & ex2 \\
4 & Ch. 3 & Data Description & ex3 \\
5 & Ch. 3 & Data Description (continued) & ex4 \\
6 & Ch. 4 & Probability and Counting Rules & ex5 \\
7 & Ch. 4 & Probability and Counting Rules (continued) & ex6 \\
8 & --- & Midterm Review & --- \\
9 & --- & Midterm Examination & --- \\
10 & Ch. 5 & Discrete Probability Distributions & ex7 \\
11 & Ch. 5 & Discrete Probability Distributions (continued) & ex8 \\
12 & Ch. 6 & The Normal Distribution & ex9 \\
13 & Ch. 6 & The Normal Distribution (continued) & ex10 \\
14 & Ch. 7 & Confidence Intervals and Sample Size & ex11 \\
15 & Ch. 7 & Confidence Intervals and Sample Size (continued) & ex12 \\
16 & --- & Final Examination & --- \\
17-18 & --- & Flexible Learning Weeks & --- \\
\end{longtable}

\subsection{Statistics (2): Chapter
Schedule}\label{statistics-2-chapter-schedule}

\begin{longtable}[]{@{}
  >{\raggedright\arraybackslash}p{(\linewidth - 6\tabcolsep) * \real{0.0594}}
  >{\raggedright\arraybackslash}p{(\linewidth - 6\tabcolsep) * \real{0.0792}}
  >{\raggedright\arraybackslash}p{(\linewidth - 6\tabcolsep) * \real{0.7624}}
  >{\raggedright\arraybackslash}p{(\linewidth - 6\tabcolsep) * \real{0.0990}}@{}}
\caption{Statistics (2) Weekly Schedule and Exercise
Mapping}\tabularnewline
\toprule\noalign{}
\begin{minipage}[b]{\linewidth}\raggedright
Week
\end{minipage} & \begin{minipage}[b]{\linewidth}\raggedright
Chapter
\end{minipage} & \begin{minipage}[b]{\linewidth}\raggedright
Topic
\end{minipage} & \begin{minipage}[b]{\linewidth}\raggedright
Exercises
\end{minipage} \\
\midrule\noalign{}
\endfirsthead
\toprule\noalign{}
\begin{minipage}[b]{\linewidth}\raggedright
Week
\end{minipage} & \begin{minipage}[b]{\linewidth}\raggedright
Chapter
\end{minipage} & \begin{minipage}[b]{\linewidth}\raggedright
Topic
\end{minipage} & \begin{minipage}[b]{\linewidth}\raggedright
Exercises
\end{minipage} \\
\midrule\noalign{}
\endhead
\bottomrule\noalign{}
\endlastfoot
1 & --- & Syllabus and Class Regulations & --- \\
2-4 & Ch. 8 & Hypothesis Testing & ex1-3 \\
5-6 & Ch. 9 & Testing the Difference Between Two Means, Two Proportions,
and Two Variances & ex4-5 \\
7 & Ch. 10 & Correlation and Regression & ex6 \\
8 & --- & Midterm Review & --- \\
9 & --- & Midterm Examination & --- \\
10-11 & Ch. 10 & Correlation and Regression (continued) & ex7-8 \\
12-13 & Ch. 11 & Other Chi-Square Tests & ex9-10 \\
14-15 & Ch. 12 & Analysis of Variance & ex11-12 \\
16 & --- & Final Examination & --- \\
17-18 & --- & Flexible Learning Weeks & --- \\
\end{longtable}

\begin{quote}
\textbf{Note:} Slide decks for Chapter 13 (Nonparametric Statistics) and
Chapter 14 (Sampling and Simulation) are maintained in the repository
and listed on the slide index as supplementary material, but are not
part of the current Statistics (2) schedule.
\end{quote}

\newpage

\section{Appendix}\label{appendix}

\subsection{A. Official Submission
Template}\label{a.-official-submission-template}

\begin{Shaded}
\begin{Highlighting}[]
\PreprocessorTok{{-}{-}{-}}
\FunctionTok{title}\KeywordTok{:}\AttributeTok{ }\StringTok{"Ex. 1"}
\FunctionTok{author}\KeywordTok{:}\AttributeTok{ }\StringTok{"A113290001"}
\FunctionTok{date}\KeywordTok{:}\AttributeTok{ }\StringTok{"2026{-}09{-}22"}
\FunctionTok{output}\KeywordTok{:}\AttributeTok{ pdf\_document}
\PreprocessorTok{{-}{-}{-}}
\end{Highlighting}
\end{Shaded}

\begin{Shaded}
\begin{Highlighting}[]
\CommentTok{\# This chunk is mandatory in every submission.}
\FunctionTok{print}\NormalTok{(uuid}\SpecialCharTok{::}\FunctionTok{UUIDgenerate}\NormalTok{())}

\CommentTok{\# Exercise solutions follow below.}
\end{Highlighting}
\end{Shaded}

\subsection{B. Directory Quick
Reference}\label{b.-directory-quick-reference}

\begin{longtable}[]{@{}
  >{\raggedright\arraybackslash}p{(\linewidth - 2\tabcolsep) * \real{0.2233}}
  >{\raggedright\arraybackslash}p{(\linewidth - 2\tabcolsep) * \real{0.7767}}@{}}
\caption{Directory Purpose Quick Reference}\tabularnewline
\toprule\noalign{}
\begin{minipage}[b]{\linewidth}\raggedright
Directory
\end{minipage} & \begin{minipage}[b]{\linewidth}\raggedright
Purpose
\end{minipage} \\
\midrule\noalign{}
\endfirsthead
\toprule\noalign{}
\begin{minipage}[b]{\linewidth}\raggedright
Directory
\end{minipage} & \begin{minipage}[b]{\linewidth}\raggedright
Purpose
\end{minipage} \\
\midrule\noalign{}
\endhead
\bottomrule\noalign{}
\endlastfoot
web/ & Website page sources (index, sylla, slide, prob, r) and RStudio
project \\
ch2/-ch14/ & Bilingual Reveal.js lecture slide sources, one subdirectory
per chapter \\
stat1-ps/ & Statistics (1) in-class exercise problem sets \\
stat2-ps/ & Statistics (2) in-class exercise problem sets \\
stat1-ex/exN/ & Statistics (1) student PDF submissions for exercise N \\
stat2-ex/exN/ & Statistics (2) student PDF submissions for exercise N \\
stat1-mid/, stat1-fin/ & Statistics (1) exam materials and grading
outputs ({[}exam{]}-score.xlsx, stat.pdf) \\
stat2-mid/, stat2-fin/ & Statistics (2) exam materials and grading
outputs ({[}exam{]}-score.xlsx, stat.pdf) \\
stu1/exN/ & Statistics (1) exercise grading outputs (scoreboard,
analytics report) \\
stu2/exN/ & Statistics (2) exercise grading outputs (scoreboard,
analytics report) \\
stu1/sem/, stu2/sem/ & Semester-aggregate scoreboard and analytics
report \\
sylla/ & Syllabus R Markdown source files \\
prompt/ & AI prompts, master slide specification, textbook PDF \\
!-Archive/ & Archive of retired materials \\
\end{longtable}

\subsection{C. References}\label{c.-references}

\begin{itemize}
\tightlist
\item
  Bluman, A. G. (2012). \emph{Elementary Statistics: A Step-by-Step
  Approach} (8th ed.). McGraw-Hill.
\item
  R Core Team. (2024). \emph{R: A language and environment for
  statistical computing}. R Foundation for Statistical Computing.
  \url{https://www.R-project.org}
\item
  Posit, PBC. (2024). \emph{RStudio Desktop}.
  \url{https://posit.co/download/rstudio-desktop}
\item
  Allaire, J. J., et al.~(2024). \emph{Quarto}. \url{https://quarto.org}
\item
  Course website: \url{https://yyliou.github.io/stat}
\end{itemize}

\end{document}
